\begin{figure}[htbp]
\centering
\begin{tikzpicture}[
  node distance=10mm and 14mm,
  box/.style={rectangle, rounded corners, draw=black!70, thick, align=center,
    text width=30mm, minimum height=13mm, font=\small\linespread{0.9}\selectfont},
  data/.style={box, fill=blue!8},
  ai/.style={box, fill=green!8},
  safe/.style={box, fill=red!8},
  outn/.style={box, fill=orange!10},
  arr/.style={-{Latex[length=2.5mm]}, thick}
]
\node[data] (knhanes) {\textbf{Stage 1}\\KNHANES\\development};
\node[data, right=of knhanes] (nhanes) {\textbf{Stage 2}\\NHANES\\external validation};
\node[data, right=of nhanes] (clinic) {\textbf{Stage 3}\\Hospital-based\\clinical cohort};

\node[ai, below=of knhanes] (harm) {Harmonized\\feature schema};
\node[ai, below=of nhanes] (model) {Interpretable +\\boosted models\\(calibrated)};
\node[ai, below=of clinic] (llm) {Open medical LLM\\schema extraction};

\node[safe, below=14mm of model] (redflag) {\textbf{Deterministic}\\\textbf{red-flag layer}\\(SSNHL / neuro)};

\node[outn, below=of harm] (assoc) {Weighted\\associations\\(no causal claim)};
\node[outn, below=of llm] (explain) {Guideline-linked\\explanations\\+ model cards};

\draw[arr] (knhanes) -- (nhanes);
\draw[arr] (nhanes) -- (clinic);
\draw[arr] (knhanes) -- (harm);
\draw[arr] (nhanes) -- (model);
\draw[arr] (clinic) -- (llm);
\draw[arr] (harm) -- (model);
\draw[arr] (llm) -- (model);
\draw[arr] (model) -- (redflag);
\draw[arr] (harm) -- (assoc);
\draw[arr] (llm) -- (explain);
\draw[arr] (redflag.west) -- ++(-1.2,0) |- (explain.east);
\end{tikzpicture}
\caption{STARS conceptual framework. Public survey data (blue) feed a harmonized schema and calibrated, interpretable models (green); an open medical language model performs schema-constrained extraction in the clinical extension. A deterministic red-flag safety layer (red) can override probabilistic output to force urgent referral, so a low predicted risk never suppresses time-critical SSNHL care. Outputs (orange) are associations and guideline-linked explanations, not diagnoses.}
\label{fig:framework}
\end{figure}

\begin{figure}[htbp]
\centering
\begin{tikzpicture}[
  every node/.style={font=\small},
  var/.style={rectangle, rounded corners, draw=black!70, thick, align=center,
    inner sep=3pt, minimum height=8mm},
  conf/.style={var, fill=blue!8},
  expo/.style={var, fill=orange!12},
  med/.style={var, fill=green!10},
  outv/.style={var, fill=red!8},
  a/.style={-{Latex[length=2mm]}, thick},
  amed/.style={-{Latex[length=2mm]}, thick, densely dashed}
]
\node[expo] (stress) at (0,0) {Perceived stress /\\work-related factors};
\node[outv] (outc) at (7.2,0) {Tinnitus \&\\hearing outcomes};
\node[med] (dep) at (3.6,-1.9) {Depressed mood, sleep\\(\emph{mediators})};
\node[conf] (conf) at (3.6,2.0) {Age, sex, education,\\occupation, noise\\(\emph{confounders})};

\draw[a] (stress) -- node[above,font=\scriptsize]{primary (direct)} (outc);
\draw[amed] (stress) -- (dep);
\draw[amed] (dep) -- (outc);
\draw[a] (conf) -- (stress);
\draw[a] (conf) -- (outc);
\end{tikzpicture}
\caption{Assumed causal directed acyclic graph (DAG). The minimal sufficient adjustment set (blue confounders) is used in the \emph{primary} analysis. Depressed mood and sleep (green) are treated as \emph{mediators} of the stress$\rightarrow$outcome path (dashed); conditioning on them (extended model) is a mediator-adjusted, potentially overadjusted estimate reported separately. Restricting to employed adults can open a selection/collider path, so all-adult analysis is primary and employed-restricted analysis is secondary.}
\label{fig:dag}
\end{figure}
